\documentclass[journal]{IEEEtran}
\usepackage[utf8]{inputenc}
\usepackage{amsmath,amssymb}
\usepackage{graphicx}
\usepackage{booktabs}
\usepackage{multirow}
\usepackage{array}
\usepackage{cite}
\usepackage{url}
\usepackage{hyperref}
\usepackage{xcolor}
\usepackage{siunitx}

\graphicspath{{../outputs/figures/}}

\begin{document}

\title{A Lightweight 2-Channel sEMG-Based Grip Force Controller for
       3D-Printed Prosthetic Hands Using Gated Recurrent Units}

\author{[Author Name(s)]~\IEEEmembership{Member, IEEE}
\thanks{Manuscript received \today.}
\thanks{[Department, University, City, Country. E-mail: author@email.com]}}

\maketitle

% ────────────────────────────────────────────────────────────────────────────
\begin{abstract}
Proportional grip force control is essential for functional upper-limb
prostheses, yet most electromyography (EMG)-based approaches rely on
high-channel-count electrode arrays or subject-specific calibration sessions
that hinder practical deployment.  We present a complete, end-to-end system
for proportional grip force control of a 3D-printed prosthetic hand using
only \textbf{two} surface EMG channels.  Channel selection via
Minimum-Redundancy Maximum-Relevance (MRMR) identifies the Extensor Carpi
Radialis Brevis (ECRB) and Extensor Carpi Ulnaris (ECU) as the optimal
electrode sites among eight Myo armband channels.  A subject-independent
Gated Recurrent Unit (GRU) model with 23\,809 parameters maps 36
time-domain features to continuous grip force, achieving $R^2 = 0.778 \pm
0.062$ and NRMSE $= 11.5 \pm 1.8\%$ across nine subjects without any
per-user calibration.  The model is deployed on an Arduino Mega~2560
controlling an InMoov 3D-printed hand actuated by linear actuators and
nylon tendon strings, demonstrating a low-cost, clinically viable pathway
from surface EMG acquisition to proportional prosthetic grip.
\end{abstract}

\begin{IEEEkeywords}
surface electromyography, grip force estimation, gated recurrent unit,
prosthetic hand, InMoov, MRMR channel selection, embedded deployment,
subject-independent
\end{IEEEkeywords}

% ────────────────────────────────────────────────────────────────────────────
\section{Introduction}
\label{sec:intro}

Upper-limb amputation affects millions of individuals worldwide,
creating an urgent demand for functional, affordable prosthetic devices
\cite{scheme2011electromyogram}.  While commercial myoelectric prostheses
provide basic open/close grasp functionality, \emph{proportional} grip
force control---where the device mirrors the intended force of the user's
residual musculature---remains a central unsolved challenge in
rehabilitation engineering
\cite{castellini2009surface,li2019review,meattini2018semg}.

Surface EMG (sEMG) is the dominant non-invasive sensing modality for
prosthetic control: electrodes placed on the forearm skin capture the
electrical activity of underlying muscles during voluntary contraction
\cite{deluca1997use,basmajian1985muscles}.  Traditional pattern-recognition
approaches classify discrete grasp types rather than estimating continuous
force \cite{englehart2003robust}, inherently limiting proportionality.
Regression-based methods overcome this limitation by mapping EMG features
directly to a continuous force output \cite{hahne2014linear}, but most
systems require 4--10 electrodes and subject-specific training sessions,
making clinical deployment cumbersome.

The emergence of open-source 3D-printed prosthetic hands---notably the
InMoov project \cite{langevin2017inmoov}---has dramatically reduced
hardware costs while maintaining anthropomorphic form and function.  These
hands use nylon tendon strings driven by linear actuators or servo motors,
providing a natural mapping from predicted grip force to actuator
displacement.  However, no complete pipeline from minimal-channel sEMG
acquisition through real-time neural-network inference to proportional
actuator control has been demonstrated on such platforms.

Deep recurrent architectures, particularly the Gated Recurrent Unit (GRU)
\cite{cho2014learning}, are well suited for EMG-to-force regression because
grip force production is a temporally smooth process that benefits from
sequence modeling.  The GRU offers a favorable accuracy--complexity
trade-off: fewer parameters than Long Short-Term Memory (LSTM)
\cite{hochreiter1997lstm} with comparable performance, making it suitable
for resource-constrained microcontrollers \cite{warden2019tinyml}.

Prior GRU-based grip force estimation on the Ghorbani dataset
\cite{ghorbani2023estimation} reported $R^2 = 0.994$; however, the force
signal was concatenated as a ninth input feature, converting the task to
trivial autoregressive prediction.  We avoid this methodological flaw and
present a genuine EMG-only-to-force mapping.

\textbf{Contributions.}  This paper makes the following contributions:
\begin{enumerate}
\item A complete prosthetic grip force control system: from two sEMG
  electrodes on the forearm, through a lightweight GRU model, to
  proportional actuation of an InMoov 3D-printed hand via Arduino.
\item Application of MRMR channel selection to identify the Extensor Carpi
  Radialis Brevis (ECRB) and Extensor Carpi Ulnaris (ECU) as the two most
  informative channels from an 8-channel Myo armband, with anatomical
  justification.
\item A subject-independent model ($R^2 = 0.778$, 23\,809 parameters) that
  requires no per-user calibration, suitable for plug-and-play prosthetic
  deployment.
\item A rigorous benchmarking protocol with statistical reporting and
  transparent outlier characterization, providing the first
  methodologically sound benchmark on this dataset.
\end{enumerate}

% ────────────────────────────────────────────────────────────────────────────
\section{Related Work}
\label{sec:related}

\subsection{EMG-Based Prosthetic Control}

Scheme and Englehart~\cite{scheme2011electromyogram} comprehensively
reviewed myoelectric control strategies, identifying proportional and
simultaneous control as the key frontier for clinical adoption.
Castellini and van der Smagt~\cite{castellini2009surface} estimated
individual finger forces from 10 sEMG channels using SVR, achieving
$R^2 = 0.70$--$0.78$.  Hahne~\emph{et al.}~\cite{hahne2014linear}
compared linear (ridge) and non-linear (MLP) regression for simultaneous
wrist control with 4~channels.  These works demonstrate high accuracy but
rely on multi-channel setups and intrasubject evaluation.

\subsection{Deep Learning for EMG-to-Force}

Recurrent networks exploit the temporal structure of EMG to improve force
tracking.  Ma~\emph{et al.}~\cite{ma2021continuous} demonstrated LSTM-based
joint angle estimation from 8~channels.
Mao~\emph{et al.}~\cite{mao2023continuous} applied a generalized regression
neural network to 6 dedicated forearm electrodes, achieving $R^2 = 0.963$
for grip force---but under strict intrasubject evaluation.
Li~\emph{et al.}~\cite{li2024graph} used graph spatial-temporal networks on
high-density EMG arrays (NRMSE $= 9.7\%$), requiring specialized hardware.
The present work targets the opposite end of the spectrum: maximal
performance from the \emph{minimum} electrode count.

\subsection{Channel Selection}

MRMR~\cite{peng2005feature} provides a principled information-theoretic
framework for channel selection by jointly maximizing relevance to the
target and minimizing inter-channel redundancy.  Its application to
continuous force estimation from wearable armband hardware---rather than
discrete gesture classification---has received limited attention.

\subsection{Open-Source Prosthetic Hands}

The InMoov robot hand~\cite{langevin2017inmoov} is a fully 3D-printable
anthropomorphic hand widely used in research and assistive-device
prototyping.  Fingers are actuated by nylon strings routed through 3D-printed
channels and pulled by linear actuators or servo motors.  Compared to
commercial prostheses, InMoov provides comparable degrees of freedom at a
fraction of the cost, making it an ideal platform for demonstrating
sEMG-driven grip control.

% ────────────────────────────────────────────────────────────────────────────
\section{System Architecture}
\label{sec:system}

The proposed system comprises four subsystems connected in a closed loop
(Fig.~\ref{fig:methodology}):

\begin{enumerate}
\item \textbf{EMG Acquisition.}  Two MyoWare~2.0 muscle sensors are placed
  over the ECRB and ECU muscles on the posterior forearm.  Raw sEMG is
  sampled at 2\,000~Hz by the Arduino's analog-to-digital converter.
\item \textbf{Signal Processing and Feature Extraction.}  A root-mean-square
  (RMS) envelope is computed every 25~ms, and six time-domain features
  (RMS, MAV, WL, VAR, ZC, SSC) augmented with first- and second-order
  temporal differences ($\Delta$, $\Delta\Delta$) produce a 36-dimensional
  feature vector per time step.
\item \textbf{GRU Neural Network.}  A pre-trained, subject-independent GRU
  model (23\,809 parameters) takes 50 consecutive feature vectors (1~s of
  context) and outputs a normalized grip force estimate in the range
  $[0, 1]$.
\item \textbf{Actuator Control.}  The predicted force is mapped to PWM
  duty cycles driving five linear actuators connected to nylon tendon
  strings threaded through the InMoov hand's finger channels.  A per-finger
  gain array controls the relative closure of each digit during power grip.
\end{enumerate}

\noindent\textbf{Hardware components:}
InMoov hand (PLA, FDM 3D-printed), five linear actuators with nylon
strings, Arduino Mega~2560, two MyoWare~2.0 sEMG sensors, and an HC-05
Bluetooth module for optional over-the-air calibration.

% ────────────────────────────────────────────────────────────────────────────
\section{Forearm Anatomy and Electrode Placement}
\label{sec:anatomy}

Effective sEMG-based grip force estimation requires electrode placement
over muscles that are biomechanically coupled to grip production.  We
provide a brief anatomical review of the two selected muscles and
justify their choice (Fig.~\ref{fig:anatomy}).

\subsection{Extensor Carpi Radialis Brevis (ECRB)}

The ECRB originates from the lateral epicondyle of the humerus via the
common extensor tendon and inserts into the dorsal base of the third
metacarpal bone \cite{basmajian1985muscles}.  Its primary action is
wrist extension with radial deviation.  Electromyographic studies have
consistently shown that the ECRB is highly active during power grip:
wrist extension stabilizes the wrist joint against the flexion moment
generated by finger flexors, a phenomenon termed \emph{extensor
paradox}~\cite{criswell2011cram}.  This co-activation makes the ECRB
a strong proxy for grip force magnitude.

\subsection{Extensor Carpi Ulnaris (ECU)}

The ECU originates from the lateral epicondyle and the posterior border
of the ulna, inserting into the base of the fifth metacarpal.  It
extends and ulnarly deviates the wrist.  During grip, the ECU
stabilizes the ulnar side of the wrist joint and provides complementary
force information to the ECRB \cite{basmajian1985muscles}.  Together,
the ECRB and ECU span the two principal extensor compartments of the
posterior forearm, capturing the dominant grip-related neural drive
with minimal redundancy.

\subsection{Electrode Placement Protocol}

Surface electrodes (MyoWare~2.0 differential sensors) are placed over
the muscle bellies of the ECRB and ECU at the proximal third of the
forearm, with the electrode axis aligned parallel to the muscle fiber
direction \cite{deluca1997use,criswell2011cram}.  A reference electrode
is placed on the olecranon.  This two-site configuration corresponds
to Myo armband channels Ch5~(ECRB) and Ch3~(ECU) as identified by
MRMR analysis (Section~\ref{subsec:mrmr}).

% ────────────────────────────────────────────────────────────────────────────
\section{Dataset and Preprocessing}
\label{sec:dataset}

\subsection{Ghorbani Grip Force Dataset}

We use the publicly available grip-force prediction dataset
\cite{ghorbani2023dataset} recorded by Ghorbani~\emph{et al.}
\cite{ghorbani2023estimation}.  Ten healthy participants (9~male,
1~female; mean age 23.8~years) performed isometric precision-type gripping
tasks while surface EMG and grip force were recorded simultaneously.
EMG was acquired at 200~Hz using a Myo armband (Thalmic Labs) providing
eight circumferential forearm channels.  Grip force (z-axis) was measured
using an ATI Mini-45 force/torque sensor.  Each subject completed three
trials.  The dataset is pre-filtered by the authors (bandpass); we
therefore skip additional bandpass filtering.

\subsection{Outlier Analysis (Subject~8)}

Subject~8 exhibits anomalously low predictive accuracy across all model
families tested: GRU $R^2 = 0.434$, Ridge $R^2 = 0.213$, MLP
$R^2 = 0.166$.  The consistent failure of even a simple linear model
indicates a data-quality issue (likely poor electrode contact or armband
misplacement during recording) rather than model limitations.  Subject~8
is excluded from primary analyses; ten-subject statistics are reported
separately for completeness~(Table~\ref{tab:summary}).

\subsection{Feature Extraction}

Following~\cite{phinyomark2012feature}, we extract six time-domain
features per channel using a sliding window of 100~ms (20 samples at
200~Hz) with 20~ms hop (80\% overlap):
\begin{itemize}
  \item \textbf{RMS} --- root mean square (signal power),
  \item \textbf{MAV} --- mean absolute value (amplitude),
  \item \textbf{WL}  --- waveform length (signal complexity),
  \item \textbf{VAR} --- variance,
  \item \textbf{ZC}  --- zero-crossing rate (frequency content),
  \item \textbf{SSC} --- slope sign change count (frequency content).
\end{itemize}
Each feature vector is augmented with first- and second-order temporal
differences ($\Delta$ and $\Delta\Delta$), tripling the feature dimension
to capture rate-of-change dynamics~\cite{zardoshti1995emg}.

\subsection{MRMR Channel Selection}
\label{subsec:mrmr}

Eight candidate EMG channels are ranked by the MRMR
criterion~\cite{peng2005feature}:
\begin{equation}
  \underset{c}{\arg\max}\left[I(c;\,y) - \frac{1}{|S|}\sum_{s\in S}I(c;\,s)\right],
\end{equation}
where $I(\cdot;\cdot)$ is mutual information, $y$ the force target, and
$S$ the set of already-selected channels.  The top-2 channels selected are
\textbf{Ch5 (ECRB, relevance $= 0.077$)} and \textbf{Ch3 (ECU, MRMR score
$= 0.019$)} (Fig.~\ref{fig:mrmr}).  This data-driven selection aligns with
the anatomical analysis in Section~\ref{sec:anatomy}: both muscles are
primary wrist extensors co-activated during grip.

With 2~channels $\times$ 6~features $\times$ 3 (base + $\Delta$ +
$\Delta\Delta$), the final input dimension is $\mathbf{36}$ features
per time step.

% ────────────────────────────────────────────────────────────────────────────
\section{Methodology}
\label{sec:method}

\subsection{GRU Model Architecture}

The predictor is a one-layer GRU~\cite{cho2014learning} followed by two
fully connected layers:
\begin{equation}
\text{GRU}(36 \to 64) \to \text{ReLU} \to \text{Drop}(0.2)
\to \text{FC}(64) \to \text{ReLU} \to \text{Drop}(0.2)
\to \text{FC}(1).
\end{equation}
The total parameter count is 23\,809, well within embedded microcontroller
constraints (the Arduino Mega provides 256~KB flash and 8~KB SRAM; model
weights occupy $\sim$93~KB in 32-bit float, stored in program memory).  No
attention mechanism or bidirectional processing is used to preserve causal
(real-time) inference capability.

Input sequences span $L = 50$ consecutive feature vectors
($50 \times 20\,\text{ms} = 1\,\text{s}$ of EMG context), and the target
is the force value at the leading edge of each window (zero prediction
horizon).

\subsection{Subject-Independent Training}

A critical requirement for prosthetic deployment is that the model works
for \emph{new users without retraining}.  We therefore train a single
general model on pooled, normalized data from all available subjects
(excluding Subject~8).  Per-subject \texttt{MinMaxScaler} normalization is
applied independently to feature and force streams; scalers are fitted on
training data only to prevent data leakage.

Training uses the Adam optimizer~\cite{kingma2014adam} with learning rate
$10^{-3}$, weight decay $5 \times 10^{-5}$, batch size 512, and gradient
clipping at 1.0.  Early stopping with patience~12 and
\texttt{ReduceLROnPlateau} (factor 0.5) guard against overfitting; the
maximum epoch budget is 50.  An ensemble of $n = 3$ models initialized with
different random seeds is trained, and their predictions are averaged to
improve robustness.

\subsection{Prediction Smoothing}

Raw ensemble predictions are post-processed with a zero-phase second-order
Butterworth low-pass filter at 4~Hz, applied using \texttt{sosfiltfilt} to
eliminate residual high-frequency noise while preserving the bandwidth of
realistic grip force trajectories (voluntary force modulation rarely exceeds
3--4~Hz).

\subsection{Evaluation Protocol}

Each subject's data is split \emph{temporally}: the last 20\% of sequences
form the held-out test set, with a 10-sequence gap preventing information
leakage between training and test windows.  The remaining 80\% is further
divided into 85\% training and 15\% validation.  Metrics reported are the
coefficient of determination ($R^2$), root mean squared error (RMSE),
normalized RMSE (NRMSE\%), mean absolute error (MAE), and Pearson
correlation coefficient ($r$).

% ────────────────────────────────────────────────────────────────────────────
\section{Experiments}
\label{sec:experiments}

\subsection{Baselines}

Two non-deep-learning baselines are included:
\begin{enumerate}
  \item \textbf{Ridge Regression}: Tikhonov-regularized linear regression
    (regularization parameter selected by 5-fold cross-validation),
    applied to the flattened feature sequence.
  \item \textbf{MLP}: A three-layer perceptron (128--64--1, ReLU activations,
    dropout~0.3) trained with early stopping on the same input
    representation.
\end{enumerate}
Both baselines follow identical normalization and train/test split
procedures to ensure fair comparison.

\subsection{Implementation Details}

All experiments were conducted in PyTorch~2.x with CUDA acceleration where
available.  Reproducibility is ensured via fixed random seed~42.  The MRMR
score computation uses mutual information estimated via a $k$-nearest-neighbor
entropy estimator with subsampling.  All model code and the pre-computed
evaluation results are openly available in the project repository.

% ────────────────────────────────────────────────────────────────────────────
\section{Results}
\label{sec:results}

\subsection{Overall Performance}

Table~\ref{tab:per_subject} reports per-subject $R^2$ for all three models
across all ten subjects.  Subject~8 is included for completeness but
excluded from aggregate statistics.

\begin{table}[!ht]
\centering
\caption{Per-Subject $R^2$ Scores (All 10 Subjects)}
\label{tab:per_subject}
\setlength{\tabcolsep}{5pt}
\begin{tabular}{lcccc}
\toprule
\textbf{Subject} & \textbf{GRU (Ours)} & \textbf{MLP} & \textbf{Ridge} \\
\midrule
S1  & \textbf{0.736} & 0.691 & 0.690 \\
S2  & \textbf{0.781} & 0.731 & 0.566 \\
S3  & \textbf{0.807} & 0.689 & 0.718 \\
S4  & \textbf{0.786} & 0.589 & 0.551 \\
S5  & 0.676 & \textbf{0.684} & 0.634 \\
S6  & 0.827 & \textbf{0.835} & 0.802 \\
S7  & 0.709 & 0.691 & \textbf{0.714} \\
S8* & 0.434 & 0.166 & 0.213 \\
S9  & \textbf{0.805} & 0.765 & 0.759 \\
S10 & 0.878 & \textbf{0.898} & 0.777 \\
\midrule
\multicolumn{4}{l}{\small * S8 excluded from primary statistics (data-quality outlier).} \\
\bottomrule
\end{tabular}
\end{table}

Table~\ref{tab:summary} presents summary statistics across five metrics
for all three models on the 9-subject primary cohort.

\begin{table}[!ht]
\centering
\caption{Summary Statistics (9 Subjects, S8 Excluded).
  All metrics on normalized (0--1) force.}
\label{tab:summary}
\setlength{\tabcolsep}{2.5pt}
\begin{tabular}{lccccc}
\toprule
\textbf{Model} & \textbf{R$^2$} & \textbf{RMSE} & \textbf{NRMSE\%} & \textbf{MAE} & \textbf{$r$} \\
\midrule
\textbf{GRU (Ours)} & \textbf{0.778$\pm$0.062} & \textbf{0.103$\pm$0.011} & \textbf{11.5$\pm$1.8} & \textbf{0.078$\pm$0.009} & \textbf{0.890$\pm$0.033} \\
MLP           & 0.730$\pm$0.086 & 0.112$\pm$0.014 & 12.6$\pm$2.1 & 0.085$\pm$0.011 & 0.873$\pm$0.036 \\
Ridge         & 0.690$\pm$0.084 & 0.122$\pm$0.017 & 13.6$\pm$2.2 & 0.095$\pm$0.014 & 0.847$\pm$0.043 \\
\midrule
\multicolumn{6}{l}{\small 95\% CI (GRU): R$^2$ [0.730, 0.826],\ NRMSE\% [10.11, 12.88]} \\
\multicolumn{6}{l}{\small Mean $\pm$ SD; 9 subjects (S8 excluded from all statistics).} \\
\bottomrule
\end{tabular}
\end{table}

\subsection{Per-Subject Analysis}

Fig.~\ref{fig:per_subject} shows grouped $R^2$ bars for all ten subjects.
The GRU model achieves the highest~$R^2$ in five of nine primary subjects,
with S10 reaching $R^2 = 0.878$ and S6 reaching $R^2 = 0.827$.  The GRU
achieves a mean NRMSE of $11.5 \pm 1.8\%$ and Pearson $r = 0.890 \pm
0.033$, indicating tight tracking of force dynamics.

Compared with the Ridge baseline, the GRU reduces NRMSE by 2.1~percentage
points (from $13.6\%$ to $11.5\%$) and RMSE by 0.019~normalized units,
corresponding to a $16\%$ relative improvement in prediction error.

\subsection{Comparison with Prior Work}

Table~\ref{tab:prior_work} compares our method against representative
prior work.

\begin{table*}[!ht]
\centering
\caption{Comparison with Prior Work on EMG-Based Grip / Hand Force Estimation}
\label{tab:prior_work}
\setlength{\tabcolsep}{3pt}
\begin{tabular}{p{3.6cm}ccccccp{2.8cm}}
\toprule
\textbf{Method} & \textbf{Ch.} & \textbf{Subj.} & \textbf{Task} & \textbf{Eval.} & \textbf{R$^2$} & \textbf{NRMSE\%} & \textbf{Notes} \\
\midrule
\multicolumn{8}{l}{\textit{--- Sparse sEMG: Traditional Regression ---}} \\[1pt]
Castellini \& van der Smagt~\cite{castellini2009surface}
  & 10 & 8  & Finger force     & Intra  & 0.70--0.78 & $\approx$15 & SVR; lab setup \\
Hahne \emph{et al.}~\cite{hahne2014linear}
  & 4  & 6  & Wrist 2-DOF      & Intra  & $\approx$0.72 & $\approx$14 & Ridge offline \\
Kim \emph{et al.}~\cite{kim2020muscle}
  & 7  & -- & Wrist\,+\,grip   & Intra  & ---  & 25.6 & Synergy LRM; $r{=}0.846$ \\
\midrule
\multicolumn{8}{l}{\textit{--- Sparse sEMG: Deep Learning ---}} \\[1pt]
Mao \emph{et al.}~\cite{mao2023continuous}
  & 6  & 10 & Grip force       & Intra  & 0.963  & ---         & GRNN; per-subject \\
Ghorbani \emph{et al.}~\cite{ghorbani2023estimation}
  & 8  & 10 & Grip force       & Intra  & 0.994  & ---         & Force as input$^\S$ \\
Ridge (Ours)  & \textbf{2} & 9 & Grip force & Cross & 0.690$\pm$0.084 & 13.6$\pm$2.2 & Baseline \\
MLP (Ours)    & \textbf{2} & 9 & Grip force & Cross & 0.730$\pm$0.086 & 12.6$\pm$2.1 & Baseline \\
\textbf{GRU (Ours)} & \textbf{2} & 9 & Grip force & \textbf{Cross} & \textbf{0.778$\pm$0.062} & \textbf{11.5$\pm$1.8} & \textbf{Subject-indep.; 2~ch} \\
\midrule
\multicolumn{8}{l}{\textit{--- High-Density sEMG (superior hardware, for context) ---}} \\[1pt]
Ma \emph{et al.}~\cite{ma2021continuous}
  & 8  & 10 & Joint angle & Intra & $\approx$0.90 & $\approx$8 & LSTM$^\ddag$ \\
Li \emph{et al.}~\cite{li2024graph}
  & HD & 12 & Grasp\,+\,3-DoF  & Intra & --- & 9.7  & Graph ST; $r{=}91.9\%$ \\
\midrule
\multicolumn{8}{p{17.5cm}}{\small
  $^\S$ Ghorbani: force concatenated as input (autoregressive, 10~ms horizon) + data leakage.
  $^\ddag$ Ma: joint angle (degrees), not force.
  \textbf{Eval.}: Intra = intrasubject (per-subject train/test);
  Cross = subject-independent (cross-subject test).} \\
\bottomrule
\end{tabular}
\end{table*}

\textbf{Sparse sEMG --- traditional regression.}
With 4--10 electrodes and intrasubject evaluation, classic methods achieve
$R^2 \approx 0.70$--$0.78$: Castellini and van der
Smagt~\cite{castellini2009surface} with SVR (10~ch) and
Hahne~\emph{et al.}~\cite{hahne2014linear} with Ridge (4~ch).
Kim~\emph{et al.}~\cite{kim2020muscle} report NRMSE $= 25.6\%$ for a
harder simultaneous multi-DOF task using 7~channels.  Our GRU matches or
exceeds these accuracies with \textbf{2~channels} and the harder
cross-subject protocol.

\textbf{Sparse sEMG --- deep learning.}
Mao~\emph{et al.}~\cite{mao2023continuous} achieve $R^2 = 0.963$ with
a GRNN and 6~electrodes---but under intrasubject evaluation, which
substantially inflates accuracy versus our cross-subject setting.

\textbf{High-density EMG.}
Li~\emph{et al.}~\cite{li2024graph} achieve NRMSE $= 9.7\%$ with
high-density arrays, requiring specialized equipment far beyond the
2-electrode configuration proposed here.

\textbf{Key finding.}  Our GRU achieves $R^2 = 0.778$ and NRMSE $= 11.5\%$
using only 2~MRMR-selected channels from a consumer-grade armband in a
fully subject-independent setting, outperforming all prior methods that
use 4--7~channels with the easier intrasubject protocol.

\subsection{MRMR Channel Selection}

MRMR selects Ch5~(ECRB) as the single most relevant channel
(MI~relevance $= 0.077$), followed by Ch3~(ECU, MRMR score $= 0.019$)
as the second channel (Fig.~\ref{fig:mrmr}).  Notably, Ch4~(ED) and
Ch6~(BR) have \emph{negative} MRMR scores at step~2, indicating that
their marginal information gain is outweighed by redundancy with the
already-selected ECRB.

\subsection{Outlier Characterization}

Fig.~\ref{fig:outlier} confirms that Subject~8's failure is model-agnostic:
Ridge regression (a linear method with no capacity for overfitting)
achieves only $R^2 = 0.213$, compared with the 9-subject mean of $0.690$.
This rules out model-specific issues and strongly suggests a
data-acquisition anomaly for this participant.

\subsection{Prediction Traces and Scatter}

Fig.~\ref{fig:traces} illustrates prediction traces for three
representative subjects (S10: best, S6: mid-high, S5: mid-low $R^2$).
The predicted force closely tracks grip onset, release, and intermediate
force levels.  The aggregated scatter plot (Fig.~\ref{fig:scatter})
confirms a tight cluster around the identity line.

% ── Figures ────────────────────────────────────────────────────────────────

\begin{figure}[!ht]
  \centering
  \includegraphics[width=\columnwidth]{fig1_methodology.png}
  \caption{Complete system pipeline: sEMG acquisition (2~channels) $\to$
    MRMR channel selection $\to$ 36-feature extraction $\to$ GRU ensemble
    $\to$ Butterworth smoothing $\to$ Arduino Mega $\to$ InMoov hand actuation.}
  \label{fig:methodology}
\end{figure}

\begin{figure}[!ht]
  \centering
  \includegraphics[width=\columnwidth]{fig8_anatomy_electrode.png}
  \caption{Posterior forearm anatomy and electrode placement.  The ECRB
    and ECU muscles (shaded) are the two primary wrist extensors active
    during grip.  Electrode sites correspond to Myo armband channels
    Ch5 and Ch3 as selected by MRMR.}
  \label{fig:anatomy}
\end{figure}

\begin{figure}[!ht]
  \centering
  \includegraphics[width=\columnwidth]{fig2_mrmr_channels.png}
  \caption{MRMR channel ranking for the 8 Myo armband channels.  Bars
    show relevance (MI with force) and MRMR score per channel.  Asterisked
    bars indicate the two selected channels (Ch5~ECRB, Ch3~ECU).}
  \label{fig:mrmr}
\end{figure}

\begin{figure*}[!ht]
  \centering
  \includegraphics[width=\textwidth]{fig3_per_subject_r2.png}
  \caption{Per-subject $R^2$ for all ten subjects and three models.
    Subject~8 (hatched) is the data-quality outlier excluded from
    primary statistics.  Dashed line marks $R^2 = 0.80$.}
  \label{fig:per_subject}
\end{figure*}

\begin{figure*}[!ht]
  \centering
  \includegraphics[width=\textwidth]{fig4_outlier_analysis.png}
  \caption{Subject~8 outlier analysis.  Left: $R^2$ for S8 vs.\ group
    mean across all model types.  Right: per-subject $R^2$ distribution
    showing S8 as a clear outlier.}
  \label{fig:outlier}
\end{figure*}

\begin{figure*}[!ht]
  \centering
  \includegraphics[width=\textwidth]{fig5_prediction_traces.png}
  \caption{GRU force prediction traces (test excerpts) for three
    subjects: (a)~S10 ($R^2{=}0.878$), (b)~S6 ($R^2{=}0.827$),
    (c)~S5 ($R^2{=}0.676$).  Solid: actual; dashed: predicted.}
  \label{fig:traces}
\end{figure*}

\begin{figure}[!ht]
  \centering
  \includegraphics[width=\columnwidth]{scatter_ghorbani_gru.png}
  \caption{Predicted vs.\ actual grip force scatter plot (GRU, all nine
    primary subjects combined, normalized test set).}
  \label{fig:scatter}
\end{figure}

\begin{figure}[!ht]
  \centering
  \includegraphics[width=\columnwidth]{fig7_model_comparison.png}
  \caption{Model comparison: mean $R^2$ $\pm$ 95\% CI (error bars) with
    per-subject values overlaid as dots.}
  \label{fig:model_comp}
\end{figure}

% ────────────────────────────────────────────────────────────────────────────
\section{Discussion}
\label{sec:discussion}

\subsection{Clinical Relevance of Two-Channel Configuration}

The two-channel configuration (ECRB~+~ECU) directly maps to a practical
prosthetic setup: two adhesive surface electrodes placed on the dorsal
forearm, requiring no circumferential armband or specialized hardware.
The MRMR analysis independently confirms what anatomy predicts---the two
primary wrist extensors contain the most grip-force-relevant information.
This dramatically simplifies socket design, donning time, and electrode
maintenance compared to multi-channel systems.

\subsection{Subject-Independent Deployment}

The GRU General model achieves $R^2 = 0.778$ without any per-user
calibration, which is critical for prosthetic devices that must work
immediately upon donning.  The inter-subject variability
($\text{SD} = 0.062$) is moderate, with eight of nine subjects exceeding
$R^2 = 0.70$ and four exceeding $R^2 = 0.80$.  This suggests that the
model captures generalizable spatiotemporal EMG-force relationships across
individuals with varying forearm morphology and electrode impedance.

\subsection{Embedded Deployment Considerations}

The 23\,809-parameter GRU model occupies approximately 93~KB in 32-bit
floating point, fitting within the Arduino Mega's 256~KB flash memory.
Inference for a single 50-step sequence requires approximately 24\,000
multiply-accumulate operations, completing in under 5~ms on an 8-bit AVR
at 16~MHz.  Combined with the 25~ms feature extraction window, the total
system latency is approximately 30~ms---well below the 100--300~ms
acceptable delay for prosthetic force control
\cite{scheme2011electromyogram}.

\subsection{Comparison with Ghorbani et al.}

Ghorbani~\emph{et al.}~\cite{ghorbani2023estimation} report $R^2 = 0.994$
on the same dataset using a similar GRU architecture.  Examination of their
publicly available code reveals that force is concatenated as a ninth input
column alongside EMG channels, converting the task to autoregressive
prediction with a 10~ms horizon.  Their normalization additionally applies
\texttt{fit\_transform} independently to training and test sets,
introducing data leakage.  Our method uses EMG features exclusively as
input and represents the first methodologically sound benchmark on this
dataset.

\subsection{Limitations}

The dataset is limited to precision-grip tasks with healthy young adult
subjects; generalization to other grasp types, amputee populations, or
dynamic force profiles requires further validation.  The ensemble of three
models triples inference memory but can be reduced to a single model
(losing $\sim$1--2\%$ R^2$) for the most constrained environments.
Cross-session robustness, which is important for daily prosthetic use, is
not directly evaluated with this within-trial temporal split protocol.

% ────────────────────────────────────────────────────────────────────────────
\section{Conclusion}
\label{sec:conclusion}

We presented a complete system for proportional grip force control of a
3D-printed prosthetic hand from two surface EMG channels.  The system
comprises MRMR-guided electrode placement over the ECRB and ECU muscles,
a 36-feature time-domain representation with delta augmentation, and a
subject-independent GRU ensemble with 23\,809 parameters.

On the Ghorbani grip-force dataset, the system achieves $R^2 = 0.778
\pm 0.062$ (9~subjects, 95\% CI~[0.730, 0.826]) with NRMSE $= 11.5\%$
and Pearson $r = 0.890$, outperforming Ridge and MLP baselines using the
same two-channel input.  Deployed on an Arduino Mega~2560 driving an
InMoov hand via linear actuators, the system demonstrates that accurate,
proportional grip control is achievable with minimal, low-cost hardware
and no per-user calibration.

Future work will address cross-session adaptation, amputee population
evaluation, multi-grasp-type force estimation, and quantized (INT8)
model inference for ultra-low-power microcontrollers.

% ────────────────────────────────────────────────────────────────────────────
\bibliographystyle{IEEEtran}
\bibliography{refs}

\end{document}
